% AML (CSAI2017P) --- Theory Quiz 4 (T04) --- 8 sets, A--H, four pages each.
% Generated from the item store; do not edit this file directly.
% Released after all sittings. Worked solutions: T04-Solutions.pdf in this folder.
\documentclass[10pt,a4paper]{article}
\usepackage[margin=16mm,top=12mm,bottom=14mm]{geometry}
\usepackage[T1]{fontenc}
\usepackage{lmodern}
\usepackage{enumitem}
\usepackage{fancyvrb}
\usepackage{amsmath}
\usepackage{amssymb}
\usepackage{array}
\usepackage{fancyhdr}
\usepackage{microtype}

\setlength{\parindent}{0pt}
\setlength{\emergencystretch}{3em}
\newcommand{\BL}{\rule[-1.5pt]{22mm}{0.4pt}}
\newcommand{\blk}[1]{\rule[-1.5pt]{#1}{0.4pt}}

\pagestyle{fancy}\fancyhf{}
\renewcommand{\headrulewidth}{0.4pt}
\renewcommand{\footrulewidth}{0pt}
\newcommand{\SetLetter}{}
\fancyhead[L]{\footnotesize CSAI2017P --- Applied Machine Learning (Theory) --- Theory Quiz 4}
\fancyhead[R]{\footnotesize\bfseries Set \SetLetter}
%% No footer. The paper is printed two A4 pages to a side of A3, four to the
%% sheet, so one A3 sheet carries one whole set and the four pages cannot be
%% separated. Identification lives once, in the name block on page 1; a
%% per-page Name/SAP ID footer and a "page n of 4" rule are both redundant
%% under that scheme, and the footer rule prints close to the fold.

\DefineVerbatimEnvironment{code}{Verbatim}%
  {fontsize=\small,xleftmargin=5mm,baselinestretch=0.98}

\newlist{qA}{enumerate}{1}
\setlist[qA]{label=\textbf{A\arabic*.},leftmargin=8mm,topsep=4pt,itemsep=9pt,parsep=1pt}
\newlist{qB}{enumerate}{1}
\setlist[qB]{label=\textbf{B\arabic*.},leftmargin=8mm,topsep=5pt,itemsep=13pt,parsep=1pt}
\newlist{qC}{enumerate}{1}
\setlist[qC]{label=\textbf{C\arabic*.},leftmargin=8mm,topsep=5pt,itemsep=8pt,parsep=1pt}
\newlist{qD}{enumerate}{1}
\setlist[qD]{label=\textbf{D\arabic*.},leftmargin=8mm,topsep=3pt,itemsep=4pt,parsep=1pt}
\newlist{qE}{enumerate}{1}
\setlist[qE]{label=\textbf{E\arabic*.},leftmargin=8mm,topsep=3pt,itemsep=5pt,parsep=1pt}

\newcommand{\parthead}[2]{\vspace{2pt}\noindent\rule{\linewidth}{0.4pt}\par\vspace{2pt}%
  \noindent{\bfseries Part #1}\hfill{\small #2}\par\vspace{1pt}}

\makeatletter
\newcommand{\rules}[2]{\par\vspace{2mm}%
  \@tempcnta=0
  \loop\ifnum\@tempcnta<#1\relax
    \noindent\rule{\linewidth}{0.3pt}\par\vspace{#2}%
    \advance\@tempcnta by 1
  \repeat}
\makeatother

\begin{document}


\renewcommand{\SetLetter}{A}\setcounter{page}{1}


\begin{center}
{\large\bfseries CSAI2017P --- Applied Machine Learning (Theory)}\\[2pt]
{\large\bfseries Theory Quiz 4 (T04) \quad\textbullet\quad Set A}\\[3pt]
{\small Duration \textbf{40 minutes} \quad\textbullet\quad Maximum marks \textbf{20}
 \quad\textbullet\quad 15 questions on four pages}
\end{center}
\vspace{2pt}

% Column widths differ slightly from LQ01's: at LQ01's widths the row runs
% ~11.6pt past the margin and "Signature" does not fit its own cell.
\noindent\begin{tabular}{@{}p{1.4cm}p{5.0cm}p{1.9cm}p{3.5cm}p{1.3cm}p{2.4cm}@{}}
\textbf{Name} & \hrulefill & \textbf{SAP ID} & \hrulefill & \textbf{Batch} & \hrulefill\\[9pt]
\textbf{Date} & \hrulefill & \textbf{Signature} & \hrulefill & \textbf{Set} & \fbox{\rule{0pt}{11pt}\ A\ }\\
\end{tabular}

\vspace{4pt}
\noindent\fbox{\parbox{\dimexpr\linewidth-2\fboxsep-2\fboxrule}{\small
\textbf{Syllabus.} The \textbf{Mid-Semester Study Material} (Lectures 1--9). Principally its later
topics: missing data, outliers and transformations, feature scaling, encoding, PCA, splitting and
cross-validation, imbalanced data, data leakage and the order of the pipeline, regression foundations,
least squares, maximum likelihood, $R^2$ and model adequacy. Lectures 1--4 are examined as retained
material.\\[3pt]
\textbf{Instructions.} Closed book. Write your answers on this paper, in the space provided. A
calculator is not needed --- the arithmetic is deliberately short. Answer every question --- there
is no negative marking. In Part C a reason is \textbf{optional}, but a correct reason can rescue a
wrong answer. In Parts D and E the marks are for the reasoning, not the vocabulary; an answer that
states what it cannot determine earns credit, one that asserts without a reason does not.\\[3pt]
\textbf{About the set letter.} Sets are handed out \textbf{at random}. Every set covers the same
topics at the same level, for the same marks, against the same marking scheme.}}
\vspace{3pt}


\parthead{A --- Multiple choice}{5 questions $\times$ 1 = 5 marks. Circle one letter.}

\begin{qA}
  \item A column has missing entries and is to be median-imputed. The median used must be computed from\\
    (a) all rows, before the split \quad (b) the training rows only\\
    (c) the test rows only \quad (d) the whole column, after the split
  \item Of the following, the method whose result is \emph{unchanged} by rescaling every feature is\\
    (a) $k$-nearest neighbours \quad (b) gradient descent on a linear model\\
    (c) a decision tree \quad (d) principal component analysis
  \item In PCA, the first principal component is the direction along which the data have\\
    (a) the largest variance \quad (b) the smallest variance\\
    (c) the largest mean \quad (d) zero correlation with $y$
  \item A training run's loss rises and becomes \texttt{nan} within a few iterations. The first thing to suspect is\\
    (a) the batch size is too small \quad (b) the model has too few parameters\\
    (c) the data are imbalanced \quad (d) $\eta$ is past the stability threshold $2/L$
  \item Grouping unlabelled customers by purchase behaviour, with no target column anywhere, is\\
    (a) supervised regression \quad (b) unsupervised learning\\
    (c) supervised classification \quad (d) reinforcement learning
\end{qA}

\vspace{4pt}

\parthead{B --- Fill in the blanks}{4 $\times$ 1 = 4 marks. Write the missing word on the rule.}

\begin{qB}

  \item For $\hat y = b_0 + b_1 x$, the least-squares slope is $b_1 = S_{xy}/\BL$.

  \item Because of the normal equations, $\mathrm{SST} = \mathrm{SSR} + \BL$ exactly.

  \item On a dataset where $95\%$ of rows are negative, a model that always predicts ``negative'' has accuracy \BL.

  \item The three families of machine learning are supervised, unsupervised and \BL.

\end{qB}

\vfill
\begin{center}\small
\begin{tabular}{|>{\centering\arraybackslash}m{1.9cm}|*{5}{>{\centering\arraybackslash}m{1.6cm}|}}
\hline
\rule{0pt}{12pt}\textbf{Question} & A1 & A2 & A3 & A4 & A5\\[2pt]\hline
\rule{0pt}{20pt}\textbf{Answer} & & & & & \\[7pt]\hline
\end{tabular}\\[2pt]
{\footnotesize Copy your circled Part A letters into this grid. The grid is what is marked.}
\end{center}

\newpage

\parthead{C --- True or false}{3 $\times$ 1 = 3 marks. Circle one. A reason is optional, but write one --- it can rescue a wrong answer.}

\begin{qC}

  \item Putting every fitted step inside a \texttt{Pipeline} removes every kind of leakage from the analysis. \hfill \textbf{TRUE} \; / \; \textbf{FALSE}\\[1pt]
    {\small Reason:}\,\blk{140mm}\rules{1}{6.4mm}

  \item $R^2$ computed on held-out data cannot be negative. \hfill \textbf{TRUE} \; / \; \textbf{FALSE}\\[1pt]
    {\small Reason:}\,\blk{140mm}\rules{1}{6.4mm}

  \item Accuracy cannot be used as the quantity a gradient-based optimiser minimises. \hfill \textbf{TRUE} \; / \; \textbf{FALSE}\\[1pt]
    {\small Reason:}\,\blk{140mm}\rules{1}{6.4mm}

\end{qC}

\vspace{4pt}

\parthead{D --- Short answer}{2 questions $\times$ 2 = 4 marks. Show how you decide, not just what you decide.}

\begin{qD}

  \item Fit $\hat y = b_0 + b_1 x$ by least squares to the four points $(1,6)$, $(3,10)$, $(5,16)$, $(7,24)$, using $b_1 = S_{xy}/S_{xx}$ and $b_0 = \bar y - b_1 \bar x$.
    \begin{enumerate}[label=(\alph*),leftmargin=7mm,itemsep=1pt,topsep=2pt]
      \item Compute $b_1$ and $b_0$, showing $S_{xx}$ and $S_{xy}$.
      \item Without working out every residual, state what $\sum_i e_i$ must come to, and say why.
    \end{enumerate}\par\vspace{62mm}

\end{qD}

\newpage

\parthead{D --- Short answer, continued}{}

\begin{qD}[start=2]

  \item A training loss curve falls quickly, then bounces up and down about a level without ever settling.
    \begin{enumerate}[label=(\alph*),leftmargin=7mm,itemsep=1pt,topsep=2pt]
      \item Give \emph{two} distinct causes that would both produce this.
      \item State one thing you could change that would tell the two apart, and say what each outcome would mean.
    \end{enumerate}\par\vspace{70mm}

\end{qD}

\vspace{4pt}

\parthead{E --- Reasoning}{1 question $\times$ 4 = 4 marks. You are marked on the reasoning you make visible.}

\begin{qE}

  \item The steps below belong to one analysis, listed in no particular order. One of the features,
\texttt{price\_per\_sqft}, was computed as $\text{Price}/\text{Area}$ --- and $\text{Price}$ is the target.

\vspace{2mm}
\hspace*{4mm}\textbf{(1)} Standardise the numeric columns. \quad \textbf{(2)} Split into training and test sets.\\
\hspace*{4mm}\textbf{(3)} Median-impute the missing areas. \quad \textbf{(4)} Drop the \texttt{ID} column.\\
\hspace*{4mm}\textbf{(5)} Oversample the minority class. \quad \textbf{(6)} Fit the model.


    \vspace{2mm}
    Answer all four parts overleaf.
    \begin{enumerate}[label=(\roman*),leftmargin=7mm,itemsep=2pt,topsep=2pt]
      \item Put the six steps into a safe order. Numbers alone are enough. \hfill \textbf{[1]}
      \item One step, placed wrongly, contaminates every fitted step after it. Name it and say why.
            \hfill \textbf{[1]}
      \item One of the problems here would survive even a perfectly built \texttt{Pipeline} with the
            steps in your order. Which, and why? \hfill \textbf{[1]}
      \item One step on the list is safe to carry out \emph{before} the split. Which, and what makes
            it different from the others? \hfill \textbf{[1]}
    \end{enumerate}

\end{qE}

\newpage

\parthead{E --- Reasoning, answer space}{Label each part (i), (ii), (iii). A sketch, a table or a list is as acceptable as prose.}

\vfill

{\footnotesize\itshape Continue on a continuation sheet if you need more room; label anything you write there with its question number.}

\newpage

\renewcommand{\SetLetter}{B}\setcounter{page}{1}


\begin{center}
{\large\bfseries CSAI2017P --- Applied Machine Learning (Theory)}\\[2pt]
{\large\bfseries Theory Quiz 4 (T04) \quad\textbullet\quad Set B}\\[3pt]
{\small Duration \textbf{40 minutes} \quad\textbullet\quad Maximum marks \textbf{20}
 \quad\textbullet\quad 15 questions on four pages}
\end{center}
\vspace{2pt}

% Column widths differ slightly from LQ01's: at LQ01's widths the row runs
% ~11.6pt past the margin and "Signature" does not fit its own cell.
\noindent\begin{tabular}{@{}p{1.4cm}p{5.0cm}p{1.9cm}p{3.5cm}p{1.3cm}p{2.4cm}@{}}
\textbf{Name} & \hrulefill & \textbf{SAP ID} & \hrulefill & \textbf{Batch} & \hrulefill\\[9pt]
\textbf{Date} & \hrulefill & \textbf{Signature} & \hrulefill & \textbf{Set} & \fbox{\rule{0pt}{11pt}\ B\ }\\
\end{tabular}

\vspace{4pt}
\noindent\fbox{\parbox{\dimexpr\linewidth-2\fboxsep-2\fboxrule}{\small
\textbf{Syllabus.} The \textbf{Mid-Semester Study Material} (Lectures 1--9). Principally its later
topics: missing data, outliers and transformations, feature scaling, encoding, PCA, splitting and
cross-validation, imbalanced data, data leakage and the order of the pipeline, regression foundations,
least squares, maximum likelihood, $R^2$ and model adequacy. Lectures 1--4 are examined as retained
material.\\[3pt]
\textbf{Instructions.} Closed book. Write your answers on this paper, in the space provided. A
calculator is not needed --- the arithmetic is deliberately short. Answer every question --- there
is no negative marking. In Part C a reason is \textbf{optional}, but a correct reason can rescue a
wrong answer. In Parts D and E the marks are for the reasoning, not the vocabulary; an answer that
states what it cannot determine earns credit, one that asserts without a reason does not.\\[3pt]
\textbf{About the set letter.} Sets are handed out \textbf{at random}. Every set covers the same
topics at the same level, for the same marks, against the same marking scheme.}}
\vspace{3pt}


\parthead{A --- Multiple choice}{5 questions $\times$ 1 = 5 marks. Circle one letter.}

\begin{qA}
  \item A table of $500$ rows is split $400/100$. A column is then median-imputed. The median must come from\\
    (a) all 500 rows \quad (b) the 100 test rows\\
    (c) the 400 training rows \quad (d) whichever 250 rows are complete
  \item A table has \texttt{Area} in the thousands and \texttt{Bedrooms} in single digits. The method least disturbed by that difference is\\
    (a) $k$-nearest neighbours \quad (b) gradient descent on a linear model\\
    (c) PCA \quad (d) a decision tree
  \item A PCA table reports explained-variance ratios $0.71,\ 0.18,\ 0.07,\ 0.04$. The first two components together carry\\
    (a) $71\%$ of the variance \quad (b) $89\%$ of the variance\\
    (c) $29\%$ of the variance \quad (d) two of the four features
  \item A log of losses reads $0.82,\ 3.9,\ 41,\ \texttt{nan}$. The first thing to suspect is\\
    (a) $\eta$ is past the stability threshold $2/L$ \quad (b) too small a batch size\\
    (c) too few parameters \quad (d) an imbalanced target
  \item A table has columns \emph{floor area}, \emph{bedrooms} and \emph{sale price}, and the task is to predict the third. This is\\
    (a) supervised classification \quad (b) unsupervised learning\\
    (c) supervised regression \quad (d) reinforcement learning
\end{qA}

\vspace{4pt}

\parthead{B --- Fill in the blanks}{4 $\times$ 1 = 4 marks. Write the missing word on the rule.}

\begin{qB}

  \item A table gives $S_{xx} = 700000$ and $S_{xy} = 42000$; the least-squares slope is $b_1 = S_{xy}/\BL = 0.06$.

  \item A table reports $\mathrm{SST} = 2626$ and $\mathrm{SSR} = 2520$, so $\mathrm{SSE} = 106$ --- because $\mathrm{SST} = \mathrm{SSR} + \BL$.

  \item A table of $1000$ patients contains $50$ who are ill. Always predicting ``not ill'' gives accuracy \BL.

  \item A table of the three families of machine learning lists supervised, unsupervised and \BL.

\end{qB}

\vfill
\begin{center}\small
\begin{tabular}{|>{\centering\arraybackslash}m{1.9cm}|*{5}{>{\centering\arraybackslash}m{1.6cm}|}}
\hline
\rule{0pt}{12pt}\textbf{Question} & A1 & A2 & A3 & A4 & A5\\[2pt]\hline
\rule{0pt}{20pt}\textbf{Answer} & & & & & \\[7pt]\hline
\end{tabular}\\[2pt]
{\footnotesize Copy your circled Part A letters into this grid. The grid is what is marked.}
\end{center}

\newpage

\parthead{C --- True or false}{3 $\times$ 1 = 3 marks. Circle one. A reason is optional, but write one --- it can rescue a wrong answer.}

\begin{qC}

  \item A results table produced entirely inside a \texttt{Pipeline} is free of every kind of leakage. \hfill \textbf{TRUE} \; / \; \textbf{FALSE}\\[1pt]
    {\small Reason:}\,\blk{140mm}\rules{1}{6.4mm}

  \item A table of held-out $R^2$ values could never legitimately contain a negative entry. \hfill \textbf{TRUE} \; / \; \textbf{FALSE}\\[1pt]
    {\small Reason:}\,\blk{140mm}\rules{1}{6.4mm}

  \item A column of accuracies cannot serve as the quantity a gradient-based optimiser minimises. \hfill \textbf{TRUE} \; / \; \textbf{FALSE}\\[1pt]
    {\small Reason:}\,\blk{140mm}\rules{1}{6.4mm}

\end{qC}

\vspace{4pt}

\parthead{D --- Short answer}{2 questions $\times$ 2 = 4 marks. Show how you decide, not just what you decide.}

\begin{qD}

  \item The table below is to be completed. Fit $\hat y = b_0 + b_1 x$ by least squares, using $b_1 = S_{xy}/S_{xx}$ and $b_0 = \bar y - b_1\bar x$.\par\vspace{1mm}\hspace*{6mm}\begin{tabular}{@{}lrrrr|r@{}}\hline\rule{0pt}{10pt}$x$ & 1 & 3 & 5 & 7 & $\sum$\\ $y$ & 8 & 10 & 14 & 20 & \\[1pt]\hline\end{tabular}
    \begin{enumerate}[label=(\alph*),leftmargin=7mm,itemsep=1pt,topsep=2pt]
      \item Compute $b_1$ and $b_0$, showing $S_{xx}$ and $S_{xy}$.
      \item Without working out every residual, state what $\sum_i e_i$ must come to, and say why.
    \end{enumerate}\par\vspace{62mm}

\end{qD}

\newpage

\parthead{D --- Short answer, continued}{}

\begin{qD}[start=2]

  \item A table of the training loss at each epoch reads $2.10,\ 0.88,\ 0.54,\ 0.61,\ 0.49,\ 0.63,\ 0.52,\ 0.60$ --- it falls, then moves up and down about a level without settling.
    \begin{enumerate}[label=(\alph*),leftmargin=7mm,itemsep=1pt,topsep=2pt]
      \item Give \emph{two} distinct causes that would both produce this.
      \item State one thing you could change that would tell the two apart, and say what each outcome would mean.
    \end{enumerate}\par\vspace{70mm}

\end{qD}

\vspace{4pt}

\parthead{E --- Reasoning}{1 question $\times$ 4 = 4 marks. You are marked on the reasoning you make visible.}

\begin{qE}

  \item A colleague's notebook performs the steps below, listed here in no particular order. One of the
features, \texttt{price\_per\_sqft}, was computed as $\text{Price}/\text{Area}$, and $\text{Price}$ is the
target.

\vspace{2mm}
\begin{center}\small
\begin{tabular}{@{}ll@{\qquad}ll@{}}
\hline
\rule{0pt}{10pt}(1) & Standardise the numeric columns & (2) & Split into training and test sets\\
(3) & Median-impute the missing areas & (4) & Drop the \texttt{ID} column\\
(5) & Oversample the minority class & (6) & Fit the model\\[1pt]
\hline
\end{tabular}
\end{center}


    \vspace{2mm}
    Answer all four parts overleaf.
    \begin{enumerate}[label=(\roman*),leftmargin=7mm,itemsep=2pt,topsep=2pt]
      \item Put the six steps into a safe order. Numbers alone are enough. \hfill \textbf{[1]}
      \item One step, placed wrongly, contaminates every fitted step after it. Name it and say why.
            \hfill \textbf{[1]}
      \item One of the problems here would survive even a perfectly built \texttt{Pipeline} with the
            steps in your order. Which, and why? \hfill \textbf{[1]}
      \item One step on the list is safe to carry out \emph{before} the split. Which, and what makes
            it different from the others? \hfill \textbf{[1]}
    \end{enumerate}

\end{qE}

\newpage

\parthead{E --- Reasoning, answer space}{Label each part (i), (ii), (iii). A sketch, a table or a list is as acceptable as prose.}

\vfill

{\footnotesize\itshape Continue on a continuation sheet if you need more room; label anything you write there with its question number.}

\newpage

\renewcommand{\SetLetter}{C}\setcounter{page}{1}


\begin{center}
{\large\bfseries CSAI2017P --- Applied Machine Learning (Theory)}\\[2pt]
{\large\bfseries Theory Quiz 4 (T04) \quad\textbullet\quad Set C}\\[3pt]
{\small Duration \textbf{40 minutes} \quad\textbullet\quad Maximum marks \textbf{20}
 \quad\textbullet\quad 15 questions on four pages}
\end{center}
\vspace{2pt}

% Column widths differ slightly from LQ01's: at LQ01's widths the row runs
% ~11.6pt past the margin and "Signature" does not fit its own cell.
\noindent\begin{tabular}{@{}p{1.4cm}p{5.0cm}p{1.9cm}p{3.5cm}p{1.3cm}p{2.4cm}@{}}
\textbf{Name} & \hrulefill & \textbf{SAP ID} & \hrulefill & \textbf{Batch} & \hrulefill\\[9pt]
\textbf{Date} & \hrulefill & \textbf{Signature} & \hrulefill & \textbf{Set} & \fbox{\rule{0pt}{11pt}\ C\ }\\
\end{tabular}

\vspace{4pt}
\noindent\fbox{\parbox{\dimexpr\linewidth-2\fboxsep-2\fboxrule}{\small
\textbf{Syllabus.} The \textbf{Mid-Semester Study Material} (Lectures 1--9). Principally its later
topics: missing data, outliers and transformations, feature scaling, encoding, PCA, splitting and
cross-validation, imbalanced data, data leakage and the order of the pipeline, regression foundations,
least squares, maximum likelihood, $R^2$ and model adequacy. Lectures 1--4 are examined as retained
material.\\[3pt]
\textbf{Instructions.} Closed book. Write your answers on this paper, in the space provided. A
calculator is not needed --- the arithmetic is deliberately short. Answer every question --- there
is no negative marking. In Part C a reason is \textbf{optional}, but a correct reason can rescue a
wrong answer. In Parts D and E the marks are for the reasoning, not the vocabulary; an answer that
states what it cannot determine earns credit, one that asserts without a reason does not.\\[3pt]
\textbf{About the set letter.} Sets are handed out \textbf{at random}. Every set covers the same
topics at the same level, for the same marks, against the same marking scheme.}}
\vspace{3pt}


\parthead{A --- Multiple choice}{5 questions $\times$ 1 = 5 marks. Circle one letter.}

\begin{qA}
  \item A histogram of a skewed column is to be median-imputed. The vertical line marking that median must be computed from\\
    (a) every row, before splitting \quad (b) the test rows only\\
    (c) the complete rows only \quad (d) the training rows only
  \item A scatter plot has one axis spanning $0$--$3$ and the other $0$--$40000$. The method whose result is unaffected by that is\\
    (a) a decision tree \quad (b) $k$-nearest neighbours\\
    (c) gradient descent on a linear model \quad (d) PCA
  \item A scree plot falls steeply and then flattens. The elbow tells you\\
    (a) which feature is the target \quad (b) how many rows to drop\\
    (c) how many components to keep \quad (d) the sign of the correlation
  \item A loss curve rises off the top of the axis and then stops being drawn at all. The first thing to suspect is\\
    (a) too small a batch size \quad (b) $\eta$ is past the stability threshold $2/L$\\
    (c) too few parameters \quad (d) an imbalanced target
  \item A scatter of unlabelled points is to be divided into three natural clumps, with no target anywhere. This is\\
    (a) supervised regression \quad (b) supervised classification\\
    (c) reinforcement learning \quad (d) unsupervised learning
\end{qA}

\vspace{4pt}

\parthead{B --- Fill in the blanks}{4 $\times$ 1 = 4 marks. Write the missing word on the rule.}

\begin{qB}

  \item On a scatter plot with the fitted line drawn, the slope of that line is $b_1 = S_{xy}/\BL$.

  \item A bar decomposing the total variation into an explained part and a leftover part shows $\mathrm{SST} = \mathrm{SSR} + \BL$.

  \item A bar chart shows $95$ negatives for every $5$ positives. A model that always predicts the majority scores an accuracy of \BL.

  \item A diagram of the three families of machine learning labels them supervised, unsupervised and \BL.

\end{qB}

\vfill
\begin{center}\small
\begin{tabular}{|>{\centering\arraybackslash}m{1.9cm}|*{5}{>{\centering\arraybackslash}m{1.6cm}|}}
\hline
\rule{0pt}{12pt}\textbf{Question} & A1 & A2 & A3 & A4 & A5\\[2pt]\hline
\rule{0pt}{20pt}\textbf{Answer} & & & & & \\[7pt]\hline
\end{tabular}\\[2pt]
{\footnotesize Copy your circled Part A letters into this grid. The grid is what is marked.}
\end{center}

\newpage

\parthead{C --- True or false}{3 $\times$ 1 = 3 marks. Circle one. A reason is optional, but write one --- it can rescue a wrong answer.}

\begin{qC}

  \item A pipeline diagram in which every fitted box sits after the split is, by that fact alone, free of all leakage. \hfill \textbf{TRUE} \; / \; \textbf{FALSE}\\[1pt]
    {\small Reason:}\,\blk{140mm}\rules{1}{6.4mm}

  \item A plot of held-out $R^2$ against model complexity can never dip below zero. \hfill \textbf{TRUE} \; / \; \textbf{FALSE}\\[1pt]
    {\small Reason:}\,\blk{140mm}\rules{1}{6.4mm}

  \item A plot of accuracy against a model parameter is a step function, which is why it cannot be the quantity minimised by gradient descent. \hfill \textbf{TRUE} \; / \; \textbf{FALSE}\\[1pt]
    {\small Reason:}\,\blk{140mm}\rules{1}{6.4mm}

\end{qC}

\vspace{4pt}

\parthead{D --- Short answer}{2 questions $\times$ 2 = 4 marks. Show how you decide, not just what you decide.}

\begin{qD}

  \item Four points are plotted: $(1,6)$, $(3,12)$, $(5,20)$, $(7,30)$. A straight line is to be drawn through them by least squares, using $b_1 = S_{xy}/S_{xx}$ and $b_0 = \bar y - b_1\bar x$.
    \begin{enumerate}[label=(\alph*),leftmargin=7mm,itemsep=1pt,topsep=2pt]
      \item Compute $b_1$ and $b_0$, showing $S_{xx}$ and $S_{xy}$.
      \item Without working out every residual, state what $\sum_i e_i$ must come to, and say why.
    \end{enumerate}\par\vspace{62mm}

\end{qD}

\newpage

\parthead{D --- Short answer, continued}{}

\begin{qD}[start=2]

  \item A plotted loss curve drops steeply for the first few epochs and then oscillates about a horizontal level, never settling onto it.
    \begin{enumerate}[label=(\alph*),leftmargin=7mm,itemsep=1pt,topsep=2pt]
      \item Give \emph{two} distinct causes that would both produce this.
      \item State one thing you could change that would tell the two apart, and say what each outcome would mean.
    \end{enumerate}\par\vspace{70mm}

\end{qD}

\vspace{4pt}

\parthead{E --- Reasoning}{1 question $\times$ 4 = 4 marks. You are marked on the reasoning you make visible.}

\begin{qE}

  \item A pipeline diagram has been drawn with its boxes in no particular order. One of the features,
\texttt{price\_per\_sqft}, was computed as $\text{Price}/\text{Area}$, and $\text{Price}$ is the target.

\vspace{2mm}
\begin{center}\small
\fbox{\ (1) standardise\ } \quad \fbox{\ (2) split\ } \quad \fbox{\ (3) median-impute\ }\\[3pt]
\fbox{\ (4) drop \texttt{ID}\ } \quad \fbox{\ (5) oversample\ } \quad \fbox{\ (6) fit the model\ }
\end{center}


    \vspace{2mm}
    Answer all four parts overleaf.
    \begin{enumerate}[label=(\roman*),leftmargin=7mm,itemsep=2pt,topsep=2pt]
      \item Put the six steps into a safe order. Numbers alone are enough. \hfill \textbf{[1]}
      \item One step, placed wrongly, contaminates every fitted step after it. Name it and say why.
            \hfill \textbf{[1]}
      \item One of the problems here would survive even a perfectly built \texttt{Pipeline} with the
            steps in your order. Which, and why? \hfill \textbf{[1]}
      \item One step on the list is safe to carry out \emph{before} the split. Which, and what makes
            it different from the others? \hfill \textbf{[1]}
    \end{enumerate}

\end{qE}

\newpage

\parthead{E --- Reasoning, answer space}{Label each part (i), (ii), (iii). A sketch, a table or a list is as acceptable as prose.}

\vfill

{\footnotesize\itshape Continue on a continuation sheet if you need more room; label anything you write there with its question number.}

\newpage

\renewcommand{\SetLetter}{D}\setcounter{page}{1}


\begin{center}
{\large\bfseries CSAI2017P --- Applied Machine Learning (Theory)}\\[2pt]
{\large\bfseries Theory Quiz 4 (T04) \quad\textbullet\quad Set D}\\[3pt]
{\small Duration \textbf{40 minutes} \quad\textbullet\quad Maximum marks \textbf{20}
 \quad\textbullet\quad 15 questions on four pages}
\end{center}
\vspace{2pt}

% Column widths differ slightly from LQ01's: at LQ01's widths the row runs
% ~11.6pt past the margin and "Signature" does not fit its own cell.
\noindent\begin{tabular}{@{}p{1.4cm}p{5.0cm}p{1.9cm}p{3.5cm}p{1.3cm}p{2.4cm}@{}}
\textbf{Name} & \hrulefill & \textbf{SAP ID} & \hrulefill & \textbf{Batch} & \hrulefill\\[9pt]
\textbf{Date} & \hrulefill & \textbf{Signature} & \hrulefill & \textbf{Set} & \fbox{\rule{0pt}{11pt}\ D\ }\\
\end{tabular}

\vspace{4pt}
\noindent\fbox{\parbox{\dimexpr\linewidth-2\fboxsep-2\fboxrule}{\small
\textbf{Syllabus.} The \textbf{Mid-Semester Study Material} (Lectures 1--9). Principally its later
topics: missing data, outliers and transformations, feature scaling, encoding, PCA, splitting and
cross-validation, imbalanced data, data leakage and the order of the pipeline, regression foundations,
least squares, maximum likelihood, $R^2$ and model adequacy. Lectures 1--4 are examined as retained
material.\\[3pt]
\textbf{Instructions.} Closed book. Write your answers on this paper, in the space provided. A
calculator is not needed --- the arithmetic is deliberately short. Answer every question --- there
is no negative marking. In Part C a reason is \textbf{optional}, but a correct reason can rescue a
wrong answer. In Parts D and E the marks are for the reasoning, not the vocabulary; an answer that
states what it cannot determine earns credit, one that asserts without a reason does not.\\[3pt]
\textbf{About the set letter.} Sets are handed out \textbf{at random}. Every set covers the same
topics at the same level, for the same marks, against the same marking scheme.}}
\vspace{3pt}


\parthead{A --- Multiple choice}{5 questions $\times$ 1 = 5 marks. Circle one letter.}

\begin{qA}
  \item An analyst fills the gaps in a patient's weight column with the median. To keep the evaluation honest, that median must come from\\
    (a) the training rows only \quad (b) all the rows she has\\
    (c) the test rows \quad (d) the rows with no gaps anywhere
  \item A bank's table mixes income in lakhs with number of dependants. The model whose answer does not change if income is rescaled is\\
    (a) $k$-nearest neighbours \quad (b) a decision tree\\
    (c) gradient descent on a linear model \quad (d) PCA
  \item A team wants to compress forty correlated sensor readings into a handful of columns that keep as much variation as possible. They want\\
    (a) oversampling \quad (b) one-hot encoding\\
    (c) a train/test split \quad (d) principal component analysis
  \item An engineer's run explodes to \texttt{nan} in the third iteration. Before anything else, she should suspect\\
    (a) too small a batch size \quad (b) too few parameters\\
    (c) $\eta$ is past the stability threshold $2/L$ \quad (d) an imbalanced target
  \item A shop wants its customers sorted into groups, but has never labelled a customer as belonging to any group. This is\\
    (a) unsupervised learning \quad (b) supervised regression\\
    (c) supervised classification \quad (d) reinforcement learning
\end{qA}

\vspace{4pt}

\parthead{B --- Fill in the blanks}{4 $\times$ 1 = 4 marks. Write the missing word on the rule.}

\begin{qB}

  \item An analyst fitting a straight line to house prices computes the slope as $b_1 = S_{xy}/\BL$.

  \item A report splits the total variation in price into the part the line explains and the part it misses: $\mathrm{SST} = \mathrm{SSR} + \BL$.

  \item A hospital screens for a disease that $5$ patients in every $100$ have. A model that always says ``healthy'' is right \BL\ of the time.

  \item A syllabus lists the three families of machine learning as supervised, unsupervised and \BL.

\end{qB}

\vfill
\begin{center}\small
\begin{tabular}{|>{\centering\arraybackslash}m{1.9cm}|*{5}{>{\centering\arraybackslash}m{1.6cm}|}}
\hline
\rule{0pt}{12pt}\textbf{Question} & A1 & A2 & A3 & A4 & A5\\[2pt]\hline
\rule{0pt}{20pt}\textbf{Answer} & & & & & \\[7pt]\hline
\end{tabular}\\[2pt]
{\footnotesize Copy your circled Part A letters into this grid. The grid is what is marked.}
\end{center}

\newpage

\parthead{C --- True or false}{3 $\times$ 1 = 3 marks. Circle one. A reason is optional, but write one --- it can rescue a wrong answer.}

\begin{qC}

  \item A consultant who has put every fitted step inside a \texttt{Pipeline} can assure the client there is no leakage of any kind. \hfill \textbf{TRUE} \; / \; \textbf{FALSE}\\[1pt]
    {\small Reason:}\,\blk{140mm}\rules{1}{6.4mm}

  \item A team reporting $R^2$ on a held-out set can be sure the figure will not come out negative. \hfill \textbf{TRUE} \; / \; \textbf{FALSE}\\[1pt]
    {\small Reason:}\,\blk{140mm}\rules{1}{6.4mm}

  \item An engineer cannot train a model by asking gradient descent to maximise accuracy directly. \hfill \textbf{TRUE} \; / \; \textbf{FALSE}\\[1pt]
    {\small Reason:}\,\blk{140mm}\rules{1}{6.4mm}

\end{qC}

\vspace{4pt}

\parthead{D --- Short answer}{2 questions $\times$ 2 = 4 marks. Show how you decide, not just what you decide.}

\begin{qD}

  \item A workshop records the number of faults found ($y$) against the number of hours spent inspecting ($x$) on four days: $(1,12)$, $(3,12)$, $(5,14)$, $(7,18)$. Fit $\hat y = b_0 + b_1 x$ by least squares, using $b_1 = S_{xy}/S_{xx}$ and $b_0 = \bar y - b_1\bar x$.
    \begin{enumerate}[label=(\alph*),leftmargin=7mm,itemsep=1pt,topsep=2pt]
      \item Compute $b_1$ and $b_0$, showing $S_{xx}$ and $S_{xy}$.
      \item Without working out every residual, state what $\sum_i e_i$ must come to, and say why.
    \end{enumerate}\par\vspace{62mm}

\end{qD}

\newpage

\parthead{D --- Short answer, continued}{}

\begin{qD}[start=2]

  \item An engineer reports that her training loss drops nicely for a few epochs and then simply bounces around a level, epoch after epoch, without ever settling.
    \begin{enumerate}[label=(\alph*),leftmargin=7mm,itemsep=1pt,topsep=2pt]
      \item Give \emph{two} distinct causes that would both produce this.
      \item State one thing you could change that would tell the two apart, and say what each outcome would mean.
    \end{enumerate}\par\vspace{70mm}

\end{qD}

\vspace{4pt}

\parthead{E --- Reasoning}{1 question $\times$ 4 = 4 marks. You are marked on the reasoning you make visible.}

\begin{qE}

  \item An intern hands you a notebook. It does the six things below, though not in this order, and one of the
features --- \texttt{price\_per\_sqft} --- was computed as $\text{Price}/\text{Area}$, where $\text{Price}$
is the very thing the model is meant to predict.

\vspace{2mm}
\hspace*{4mm}\textbf{(1)} Standardise the numeric columns. \quad \textbf{(2)} Split into training and test sets.\\
\hspace*{4mm}\textbf{(3)} Median-impute the missing areas. \quad \textbf{(4)} Drop the \texttt{ID} column.\\
\hspace*{4mm}\textbf{(5)} Oversample the minority class. \quad \textbf{(6)} Fit the model.


    \vspace{2mm}
    Answer all four parts overleaf.
    \begin{enumerate}[label=(\roman*),leftmargin=7mm,itemsep=2pt,topsep=2pt]
      \item Put the six steps into a safe order. Numbers alone are enough. \hfill \textbf{[1]}
      \item One step, placed wrongly, contaminates every fitted step after it. Name it and say why.
            \hfill \textbf{[1]}
      \item One of the problems here would survive even a perfectly built \texttt{Pipeline} with the
            steps in your order. Which, and why? \hfill \textbf{[1]}
      \item One step on the list is safe to carry out \emph{before} the split. Which, and what makes
            it different from the others? \hfill \textbf{[1]}
    \end{enumerate}

\end{qE}

\newpage

\parthead{E --- Reasoning, answer space}{Label each part (i), (ii), (iii). A sketch, a table or a list is as acceptable as prose.}

\vfill

{\footnotesize\itshape Continue on a continuation sheet if you need more room; label anything you write there with its question number.}

\newpage

\renewcommand{\SetLetter}{E}\setcounter{page}{1}


\begin{center}
{\large\bfseries CSAI2017P --- Applied Machine Learning (Theory)}\\[2pt]
{\large\bfseries Theory Quiz 4 (T04) \quad\textbullet\quad Set E}\\[3pt]
{\small Duration \textbf{40 minutes} \quad\textbullet\quad Maximum marks \textbf{20}
 \quad\textbullet\quad 15 questions on four pages}
\end{center}
\vspace{2pt}

% Column widths differ slightly from LQ01's: at LQ01's widths the row runs
% ~11.6pt past the margin and "Signature" does not fit its own cell.
\noindent\begin{tabular}{@{}p{1.4cm}p{5.0cm}p{1.9cm}p{3.5cm}p{1.3cm}p{2.4cm}@{}}
\textbf{Name} & \hrulefill & \textbf{SAP ID} & \hrulefill & \textbf{Batch} & \hrulefill\\[9pt]
\textbf{Date} & \hrulefill & \textbf{Signature} & \hrulefill & \textbf{Set} & \fbox{\rule{0pt}{11pt}\ E\ }\\
\end{tabular}

\vspace{4pt}
\noindent\fbox{\parbox{\dimexpr\linewidth-2\fboxsep-2\fboxrule}{\small
\textbf{Syllabus.} The \textbf{Mid-Semester Study Material} (Lectures 1--9). Principally its later
topics: missing data, outliers and transformations, feature scaling, encoding, PCA, splitting and
cross-validation, imbalanced data, data leakage and the order of the pipeline, regression foundations,
least squares, maximum likelihood, $R^2$ and model adequacy. Lectures 1--4 are examined as retained
material.\\[3pt]
\textbf{Instructions.} Closed book. Write your answers on this paper, in the space provided. A
calculator is not needed --- the arithmetic is deliberately short. Answer every question --- there
is no negative marking. In Part C a reason is \textbf{optional}, but a correct reason can rescue a
wrong answer. In Parts D and E the marks are for the reasoning, not the vocabulary; an answer that
states what it cannot determine earns credit, one that asserts without a reason does not.\\[3pt]
\textbf{About the set letter.} Sets are handed out \textbf{at random}. Every set covers the same
topics at the same level, for the same marks, against the same marking scheme.}}
\vspace{3pt}


\parthead{A --- Multiple choice}{5 questions $\times$ 1 = 5 marks. Circle one letter.}

\begin{qA}
  \item \texttt{SimpleImputer(strategy='median')} is used inside a \texttt{Pipeline}. The median it applies to the test rows is computed from\\
    (a) all rows passed to \texttt{fit\_transform} \quad (b) the training rows only\\
    (c) the test rows themselves \quad (d) the complete rows only
  \item Of \texttt{DecisionTreeRegressor}, \texttt{KNeighborsRegressor}, \texttt{SGDRegressor} and \texttt{PCA}, the one that does \emph{not} need \texttt{StandardScaler} before it is\\
    (a) \texttt{KNeighborsRegressor} \quad (b) \texttt{SGDRegressor}\\
    (c) \texttt{PCA} \quad (d) \texttt{DecisionTreeRegressor}
  \item \texttt{PCA(n\_components=2).fit(X)} then \texttt{pca.explained\_variance\_ratio\_} returns \texttt{[0.71, 0.18]}. Together the two components carry\\
    (a) $89\%$ of the variance \quad (b) $71\%$ of the variance\\
    (c) $29\%$ of the variance \quad (d) two original features
  \item A training loop prints \texttt{0.82, 3.9, 41.0, nan}. The first thing to suspect is\\
    (a) the batch size is too small \quad (b) the model has too few parameters\\
    (c) \texttt{eta} is past $2/L$ \quad (d) the classes are imbalanced
  \item \texttt{KMeans(n\_clusters=3).fit(X)} is called and no \texttt{y} exists anywhere in the script. The task is\\
    (a) supervised regression \quad (b) supervised classification\\
    (c) reinforcement learning \quad (d) unsupervised learning
\end{qA}

\vspace{4pt}

\parthead{B --- Fill in the blanks}{4 $\times$ 1 = 4 marks. Write the missing word on the rule.}

\begin{qB}

  \item With \texttt{Sxy} and \texttt{Sxx} already computed, the slope is \texttt{b1 = Sxy / }\BL.

  \item \texttt{sst == ssr + sse} evaluates to \texttt{True} for a least-squares fit, because $\mathrm{SST} = \mathrm{SSR} + \BL$.

  \item With \texttt{y.mean() == 0.05}, \texttt{DummyClassifier(strategy='most\_frequent')} scores an accuracy of \BL.

  \item Beyond supervised and unsupervised, the third family of machine learning is \BL.

\end{qB}

\vfill
\begin{center}\small
\begin{tabular}{|>{\centering\arraybackslash}m{1.9cm}|*{5}{>{\centering\arraybackslash}m{1.6cm}|}}
\hline
\rule{0pt}{12pt}\textbf{Question} & A1 & A2 & A3 & A4 & A5\\[2pt]\hline
\rule{0pt}{20pt}\textbf{Answer} & & & & & \\[7pt]\hline
\end{tabular}\\[2pt]
{\footnotesize Copy your circled Part A letters into this grid. The grid is what is marked.}
\end{center}

\newpage

\parthead{C --- True or false}{3 $\times$ 1 = 3 marks. Circle one. A reason is optional, but write one --- it can rescue a wrong answer.}

\begin{qC}

  \item Wrapping every step in \texttt{make\_pipeline(...)} guarantees that no leakage of any kind remains. \hfill \textbf{TRUE} \; / \; \textbf{FALSE}\\[1pt]
    {\small Reason:}\,\blk{140mm}\rules{1}{6.4mm}

  \item \texttt{r2\_score(y\_test, pred)} can never return a negative number. \hfill \textbf{TRUE} \; / \; \textbf{FALSE}\\[1pt]
    {\small Reason:}\,\blk{140mm}\rules{1}{6.4mm}

  \item \texttt{accuracy\_score} cannot be handed to a gradient-based optimiser as the objective to minimise. \hfill \textbf{TRUE} \; / \; \textbf{FALSE}\\[1pt]
    {\small Reason:}\,\blk{140mm}\rules{1}{6.4mm}

\end{qC}

\vspace{4pt}

\parthead{D --- Short answer}{2 questions $\times$ 2 = 4 marks. Show how you decide, not just what you decide.}

\begin{qD}

  \item \texttt{x = np.array([1, 3, 5, 7])} and \texttt{y = np.array([6, 14, 24, 36])}. Fit $\hat y = b_0 + b_1x$ by hand, using $b_1 = S_{xy}/S_{xx}$ and $b_0 = \bar y - b_1\bar x$ --- that is, what this would print:
\begin{code}
b1 = ((x - x.mean()) * (y - y.mean())).sum() / ((x - x.mean())**2).sum()
b0 = y.mean() - b1 * x.mean()
print(b1, b0)
\end{code}
    \begin{enumerate}[label=(\alph*),leftmargin=7mm,itemsep=1pt,topsep=2pt]
      \item Compute $b_1$ and $b_0$, showing $S_{xx}$ and $S_{xy}$.
      \item Without working out every residual, state what $\sum_i e_i$ must come to, and say why.
    \end{enumerate}\par\vspace{62mm}

\end{qD}

\newpage

\parthead{D --- Short answer, continued}{}

\begin{qD}[start=2]

  \item A training loop prints the loss each epoch: it falls steeply, then prints values that move up and down about one level and stay there.
    \begin{enumerate}[label=(\alph*),leftmargin=7mm,itemsep=1pt,topsep=2pt]
      \item Give \emph{two} distinct causes that would both produce this.
      \item State one thing you could change that would tell the two apart, and say what each outcome would mean.
    \end{enumerate}\par\vspace{70mm}

\end{qD}

\vspace{4pt}

\parthead{E --- Reasoning}{1 question $\times$ 4 = 4 marks. You are marked on the reasoning you make visible.}

\begin{qE}

  \item The cell below runs without error. \texttt{price\_per\_sqft} is computed from \texttt{Price\_lakh},
which is the target.

\begin{code}
df["price_per_sqft"] = df["Price_lakh"] / df["Area_sqft"]
df = df.fillna(df.median(numeric_only=True))
X = StandardScaler().fit_transform(df.drop(columns=["Price_lakh"]))
X_tr, X_te, y_tr, y_te = train_test_split(X, df["Price_lakh"])
\end{code}


    \vspace{2mm}
    Answer all four parts overleaf.
    \begin{enumerate}[label=(\roman*),leftmargin=7mm,itemsep=2pt,topsep=2pt]
      \item Put the six steps into a safe order. Numbers alone are enough. \hfill \textbf{[1]}
      \item One step, placed wrongly, contaminates every fitted step after it. Name it and say why.
            \hfill \textbf{[1]}
      \item One of the problems here would survive even a perfectly built \texttt{Pipeline} with the
            steps in your order. Which, and why? \hfill \textbf{[1]}
      \item One step on the list is safe to carry out \emph{before} the split. Which, and what makes
            it different from the others? \hfill \textbf{[1]}
    \end{enumerate}

\end{qE}

\newpage

\parthead{E --- Reasoning, answer space}{Label each part (i), (ii), (iii). A sketch, a table or a list is as acceptable as prose.}

\vfill

{\footnotesize\itshape Continue on a continuation sheet if you need more room; label anything you write there with its question number.}

\newpage

\renewcommand{\SetLetter}{F}\setcounter{page}{1}


\begin{center}
{\large\bfseries CSAI2017P --- Applied Machine Learning (Theory)}\\[2pt]
{\large\bfseries Theory Quiz 4 (T04) \quad\textbullet\quad Set F}\\[3pt]
{\small Duration \textbf{40 minutes} \quad\textbullet\quad Maximum marks \textbf{20}
 \quad\textbullet\quad 15 questions on four pages}
\end{center}
\vspace{2pt}

% Column widths differ slightly from LQ01's: at LQ01's widths the row runs
% ~11.6pt past the margin and "Signature" does not fit its own cell.
\noindent\begin{tabular}{@{}p{1.4cm}p{5.0cm}p{1.9cm}p{3.5cm}p{1.3cm}p{2.4cm}@{}}
\textbf{Name} & \hrulefill & \textbf{SAP ID} & \hrulefill & \textbf{Batch} & \hrulefill\\[9pt]
\textbf{Date} & \hrulefill & \textbf{Signature} & \hrulefill & \textbf{Set} & \fbox{\rule{0pt}{11pt}\ F\ }\\
\end{tabular}

\vspace{4pt}
\noindent\fbox{\parbox{\dimexpr\linewidth-2\fboxsep-2\fboxrule}{\small
\textbf{Syllabus.} The \textbf{Mid-Semester Study Material} (Lectures 1--9). Principally its later
topics: missing data, outliers and transformations, feature scaling, encoding, PCA, splitting and
cross-validation, imbalanced data, data leakage and the order of the pipeline, regression foundations,
least squares, maximum likelihood, $R^2$ and model adequacy. Lectures 1--4 are examined as retained
material.\\[3pt]
\textbf{Instructions.} Closed book. Write your answers on this paper, in the space provided. A
calculator is not needed --- the arithmetic is deliberately short. Answer every question --- there
is no negative marking. In Part C a reason is \textbf{optional}, but a correct reason can rescue a
wrong answer. In Parts D and E the marks are for the reasoning, not the vocabulary; an answer that
states what it cannot determine earns credit, one that asserts without a reason does not.\\[3pt]
\textbf{About the set letter.} Sets are handed out \textbf{at random}. Every set covers the same
topics at the same level, for the same marks, against the same marking scheme.}}
\vspace{3pt}


\parthead{A --- Multiple choice}{5 questions $\times$ 1 = 5 marks. Circle one letter.}

\begin{qA}
  \item Gaps in a column are to be filled with its middle value. For the test score to mean anything, that middle value must be worked out from\\
    (a) every row there is \quad (b) the held-back rows\\
    (c) the training rows alone \quad (d) only the rows with nothing missing
  \item One column counts in thousands and another in single digits. The method that is untroubled by the difference is\\
    (a) a decision tree \quad (b) nearest neighbours\\
    (c) gradient descent on a linear model \quad (d) principal components
  \item Forty columns that largely repeat one another are to be replaced by a handful that keep as much of the variation as possible. That is\\
    (a) oversampling \quad (b) principal component analysis\\
    (c) one-hot encoding \quad (d) splitting the data
  \item A run's loss grows larger every step and then stops being a number at all. The first suspect is\\
    (a) the batches are too small \quad (b) the model is too small\\
    (c) the classes are unequal \quad (d) the step size is past the stability limit
  \item Customers are to be sorted into groups that nobody has ever named or labelled. This is\\
    (a) unsupervised learning \quad (b) supervised regression\\
    (c) supervised classification \quad (d) reinforcement learning
\end{qA}

\vspace{4pt}

\parthead{B --- Fill in the blanks}{4 $\times$ 1 = 4 marks. Write the missing word on the rule.}

\begin{qB}

  \item The slope of a least-squares line is the cross-product sum divided by \BL.

  \item The total variation splits exactly into the part the line explains and the part it leaves over: $\mathrm{SST} = \mathrm{SSR} + \BL$.

  \item Where nineteen cases in twenty are negative, always answering ``negative'' is right \BL\ of the time.

  \item After supervised and unsupervised, the third family of machine learning is \BL.

\end{qB}

\vfill
\begin{center}\small
\begin{tabular}{|>{\centering\arraybackslash}m{1.9cm}|*{5}{>{\centering\arraybackslash}m{1.6cm}|}}
\hline
\rule{0pt}{12pt}\textbf{Question} & A1 & A2 & A3 & A4 & A5\\[2pt]\hline
\rule{0pt}{20pt}\textbf{Answer} & & & & & \\[7pt]\hline
\end{tabular}\\[2pt]
{\footnotesize Copy your circled Part A letters into this grid. The grid is what is marked.}
\end{center}

\newpage

\parthead{C --- True or false}{3 $\times$ 1 = 3 marks. Circle one. A reason is optional, but write one --- it can rescue a wrong answer.}

\begin{qC}

  \item Once every fitted step has been placed inside a pipeline, no kind of leakage can remain. \hfill \textbf{TRUE} \; / \; \textbf{FALSE}\\[1pt]
    {\small Reason:}\,\blk{140mm}\rules{1}{6.4mm}

  \item The share of variation explained, measured on data the model has never seen, can never come out below zero. \hfill \textbf{TRUE} \; / \; \textbf{FALSE}\\[1pt]
    {\small Reason:}\,\blk{140mm}\rules{1}{6.4mm}

  \item The proportion of cases got right cannot itself be the quantity a gradient method is asked to improve. \hfill \textbf{TRUE} \; / \; \textbf{FALSE}\\[1pt]
    {\small Reason:}\,\blk{140mm}\rules{1}{6.4mm}

\end{qC}

\vspace{4pt}

\parthead{D --- Short answer}{2 questions $\times$ 2 = 4 marks. Show how you decide, not just what you decide.}

\begin{qD}

  \item Four readings were taken. At $x = 1, 3, 5, 7$ the measured values were $7$, $1$, $-3$ and $-5$. Fit a straight line by least squares: the slope is the cross-product sum divided by the sum of squared deviations in $x$, and the intercept then follows from the line passing through the point of means.
    \begin{enumerate}[label=(\alph*),leftmargin=7mm,itemsep=1pt,topsep=2pt]
      \item Compute $b_1$ and $b_0$, showing $S_{xx}$ and $S_{xy}$.
      \item Without working out every residual, state what $\sum_i e_i$ must come to, and say why.
    \end{enumerate}\par\vspace{62mm}

\end{qD}

\newpage

\parthead{D --- Short answer, continued}{}

\begin{qD}[start=2]

  \item A training run's loss falls sharply at first and then, epoch after epoch, simply moves up and down about one level without ever settling on it.
    \begin{enumerate}[label=(\alph*),leftmargin=7mm,itemsep=1pt,topsep=2pt]
      \item Give \emph{two} distinct causes that would both produce this.
      \item State one thing you could change that would tell the two apart, and say what each outcome would mean.
    \end{enumerate}\par\vspace{70mm}

\end{qD}

\vspace{4pt}

\parthead{E --- Reasoning}{1 question $\times$ 4 = 4 marks. You are marked on the reasoning you make visible.}

\begin{qE}

  \item A colleague describes what her notebook does, though not in the order she does it. One of the columns
she feeds the model is the price divided by the area --- and the price is the very thing she is trying to
predict.

\vspace{2mm}
\hspace*{4mm}\textbf{(1)} Put every numeric column on a common scale. \quad \textbf{(2)} Set part of the data aside for testing.\\
\hspace*{4mm}\textbf{(3)} Fill the missing areas with the middle value. \quad \textbf{(4)} Throw away the identifier column.\\
\hspace*{4mm}\textbf{(5)} Duplicate rows of the rarer class. \quad \textbf{(6)} Fit the model.


    \vspace{2mm}
    Answer all four parts overleaf.
    \begin{enumerate}[label=(\roman*),leftmargin=7mm,itemsep=2pt,topsep=2pt]
      \item Put the six steps into a safe order. Numbers alone are enough. \hfill \textbf{[1]}
      \item One step, placed wrongly, contaminates every fitted step after it. Name it and say why.
            \hfill \textbf{[1]}
      \item One of the problems here would survive even a perfectly built \texttt{Pipeline} with the
            steps in your order. Which, and why? \hfill \textbf{[1]}
      \item One step on the list is safe to carry out \emph{before} the split. Which, and what makes
            it different from the others? \hfill \textbf{[1]}
    \end{enumerate}

\end{qE}

\newpage

\parthead{E --- Reasoning, answer space}{Label each part (i), (ii), (iii). A sketch, a table or a list is as acceptable as prose.}

\vfill

{\footnotesize\itshape Continue on a continuation sheet if you need more room; label anything you write there with its question number.}

\newpage

\renewcommand{\SetLetter}{G}\setcounter{page}{1}


\begin{center}
{\large\bfseries CSAI2017P --- Applied Machine Learning (Theory)}\\[2pt]
{\large\bfseries Theory Quiz 4 (T04) \quad\textbullet\quad Set G}\\[3pt]
{\small Duration \textbf{40 minutes} \quad\textbullet\quad Maximum marks \textbf{20}
 \quad\textbullet\quad 15 questions on four pages}
\end{center}
\vspace{2pt}

% Column widths differ slightly from LQ01's: at LQ01's widths the row runs
% ~11.6pt past the margin and "Signature" does not fit its own cell.
\noindent\begin{tabular}{@{}p{1.4cm}p{5.0cm}p{1.9cm}p{3.5cm}p{1.3cm}p{2.4cm}@{}}
\textbf{Name} & \hrulefill & \textbf{SAP ID} & \hrulefill & \textbf{Batch} & \hrulefill\\[9pt]
\textbf{Date} & \hrulefill & \textbf{Signature} & \hrulefill & \textbf{Set} & \fbox{\rule{0pt}{11pt}\ G\ }\\
\end{tabular}

\vspace{4pt}
\noindent\fbox{\parbox{\dimexpr\linewidth-2\fboxsep-2\fboxrule}{\small
\textbf{Syllabus.} The \textbf{Mid-Semester Study Material} (Lectures 1--9). Principally its later
topics: missing data, outliers and transformations, feature scaling, encoding, PCA, splitting and
cross-validation, imbalanced data, data leakage and the order of the pipeline, regression foundations,
least squares, maximum likelihood, $R^2$ and model adequacy. Lectures 1--4 are examined as retained
material.\\[3pt]
\textbf{Instructions.} Closed book. Write your answers on this paper, in the space provided. A
calculator is not needed --- the arithmetic is deliberately short. Answer every question --- there
is no negative marking. In Part C a reason is \textbf{optional}, but a correct reason can rescue a
wrong answer. In Parts D and E the marks are for the reasoning, not the vocabulary; an answer that
states what it cannot determine earns credit, one that asserts without a reason does not.\\[3pt]
\textbf{About the set letter.} Sets are handed out \textbf{at random}. Every set covers the same
topics at the same level, for the same marks, against the same marking scheme.}}
\vspace{3pt}


\parthead{A --- Multiple choice}{5 questions $\times$ 1 = 5 marks. Circle one letter.}

\begin{qA}
  \item Compare imputing with the median of the training rows against imputing with the median of all rows\\
    (a) only the second gives an honest test score \quad (b) both give the same score\\
    (c) neither can be done before a split \quad (d) only the first gives an honest test score
  \item Compare a decision tree with $k$-nearest neighbours on unscaled data\\
    (a) the neighbours are unaffected, the tree is not \quad (b) the tree is unaffected, the neighbours are not\\
    (c) both are unaffected \quad (d) both are equally disturbed
  \item Compare the first principal component with the second\\
    (a) the second carries more variance \quad (b) they carry equal variance\\
    (c) the first carries more variance \quad (d) the first is always a raw feature
  \item Compare a loss that flattens with one that rises to \texttt{nan}: the second means\\
    (a) $\eta$ is past $2/L$ \quad (b) the batch size is too small\\
    (c) the model is too small \quad (d) the classes are imbalanced
  \item Compare (i) sorting unlabelled customers into groups with (ii) predicting a house's sale price\\
    (a) (i) supervised, (ii) unsupervised \quad (b) (i) unsupervised, (ii) supervised regression\\
    (c) both unsupervised \quad (d) both supervised classification
\end{qA}

\vspace{4pt}

\parthead{B --- Fill in the blanks}{4 $\times$ 1 = 4 marks. Write the missing word on the rule.}

\begin{qB}

  \item Compare $S_{xy}$ with the least-squares slope: the second is the first divided by \BL.

  \item Compare $\mathrm{SST}$ with $\mathrm{SSR}$: the difference between them is exactly \BL.

  \item Compare a real model with one that always predicts the majority, where $95$ rows in $100$ are negative: the second scores \BL.

  \item Compare the three families of learning: supervised, unsupervised, and \BL.

\end{qB}

\vfill
\begin{center}\small
\begin{tabular}{|>{\centering\arraybackslash}m{1.9cm}|*{5}{>{\centering\arraybackslash}m{1.6cm}|}}
\hline
\rule{0pt}{12pt}\textbf{Question} & A1 & A2 & A3 & A4 & A5\\[2pt]\hline
\rule{0pt}{20pt}\textbf{Answer} & & & & & \\[7pt]\hline
\end{tabular}\\[2pt]
{\footnotesize Copy your circled Part A letters into this grid. The grid is what is marked.}
\end{center}

\newpage

\parthead{C --- True or false}{3 $\times$ 1 = 3 marks. Circle one. A reason is optional, but write one --- it can rescue a wrong answer.}

\begin{qC}

  \item Between target leakage and preprocessing leakage, a \texttt{Pipeline} prevents both. \hfill \textbf{TRUE} \; / \; \textbf{FALSE}\\[1pt]
    {\small Reason:}\,\blk{140mm}\rules{1}{6.4mm}

  \item Between training $R^2$ and held-out $R^2$, neither can be negative. \hfill \textbf{TRUE} \; / \; \textbf{FALSE}\\[1pt]
    {\small Reason:}\,\blk{140mm}\rules{1}{6.4mm}

  \item Between cross-entropy and accuracy, only the first can be the quantity a gradient method minimises. \hfill \textbf{TRUE} \; / \; \textbf{FALSE}\\[1pt]
    {\small Reason:}\,\blk{140mm}\rules{1}{6.4mm}

\end{qC}

\vspace{4pt}

\parthead{D --- Short answer}{2 questions $\times$ 2 = 4 marks. Show how you decide, not just what you decide.}

\begin{qD}

  \item Compare the four points $(1,10)$, $(3,20)$, $(5,32)$, $(7,46)$ with the least-squares line through them. Fit $\hat y = b_0 + b_1x$, using $b_1 = S_{xy}/S_{xx}$ and $b_0 = \bar y - b_1\bar x$.
    \begin{enumerate}[label=(\alph*),leftmargin=7mm,itemsep=1pt,topsep=2pt]
      \item Compute $b_1$ and $b_0$, showing $S_{xx}$ and $S_{xy}$.
      \item Without working out every residual, state what $\sum_i e_i$ must come to, and say why.
    \end{enumerate}\par\vspace{62mm}

\end{qD}

\newpage

\parthead{D --- Short answer, continued}{}

\begin{qD}[start=2]

  \item Compare two loss curves from the same model: one falls and flattens onto a level, the other falls and then bounces up and down about a level without settling.
    \begin{enumerate}[label=(\alph*),leftmargin=7mm,itemsep=1pt,topsep=2pt]
      \item Give \emph{two} distinct causes that would both produce this.
      \item State one thing you could change that would tell the two apart, and say what each outcome would mean.
    \end{enumerate}\par\vspace{70mm}

\end{qD}

\vspace{4pt}

\parthead{E --- Reasoning}{1 question $\times$ 4 = 4 marks. You are marked on the reasoning you make visible.}

\begin{qE}

  \item Two colleagues run the same six steps in different orders, and get very different test scores. One of
the features, \texttt{price\_per\_sqft}, was computed as $\text{Price}/\text{Area}$, and $\text{Price}$ is
the target.

\vspace{2mm}
\hspace*{4mm}\textbf{(1)} Standardise the numeric columns. \quad \textbf{(2)} Split into training and test sets.\\
\hspace*{4mm}\textbf{(3)} Median-impute the missing areas. \quad \textbf{(4)} Drop the \texttt{ID} column.\\
\hspace*{4mm}\textbf{(5)} Oversample the minority class. \quad \textbf{(6)} Fit the model.


    \vspace{2mm}
    Answer all four parts overleaf.
    \begin{enumerate}[label=(\roman*),leftmargin=7mm,itemsep=2pt,topsep=2pt]
      \item Put the six steps into a safe order. Numbers alone are enough. \hfill \textbf{[1]}
      \item One step, placed wrongly, contaminates every fitted step after it. Name it and say why.
            \hfill \textbf{[1]}
      \item One of the problems here would survive even a perfectly built \texttt{Pipeline} with the
            steps in your order. Which, and why? \hfill \textbf{[1]}
      \item One step on the list is safe to carry out \emph{before} the split. Which, and what makes
            it different from the others? \hfill \textbf{[1]}
    \end{enumerate}

\end{qE}

\newpage

\parthead{E --- Reasoning, answer space}{Label each part (i), (ii), (iii). A sketch, a table or a list is as acceptable as prose.}

\vfill

{\footnotesize\itshape Continue on a continuation sheet if you need more room; label anything you write there with its question number.}

\newpage

\renewcommand{\SetLetter}{H}\setcounter{page}{1}


\begin{center}
{\large\bfseries CSAI2017P --- Applied Machine Learning (Theory)}\\[2pt]
{\large\bfseries Theory Quiz 4 (T04) \quad\textbullet\quad Set H}\\[3pt]
{\small Duration \textbf{40 minutes} \quad\textbullet\quad Maximum marks \textbf{20}
 \quad\textbullet\quad 15 questions on four pages}
\end{center}
\vspace{2pt}

% Column widths differ slightly from LQ01's: at LQ01's widths the row runs
% ~11.6pt past the margin and "Signature" does not fit its own cell.
\noindent\begin{tabular}{@{}p{1.4cm}p{5.0cm}p{1.9cm}p{3.5cm}p{1.3cm}p{2.4cm}@{}}
\textbf{Name} & \hrulefill & \textbf{SAP ID} & \hrulefill & \textbf{Batch} & \hrulefill\\[9pt]
\textbf{Date} & \hrulefill & \textbf{Signature} & \hrulefill & \textbf{Set} & \fbox{\rule{0pt}{11pt}\ H\ }\\
\end{tabular}

\vspace{4pt}
\noindent\fbox{\parbox{\dimexpr\linewidth-2\fboxsep-2\fboxrule}{\small
\textbf{Syllabus.} The \textbf{Mid-Semester Study Material} (Lectures 1--9). Principally its later
topics: missing data, outliers and transformations, feature scaling, encoding, PCA, splitting and
cross-validation, imbalanced data, data leakage and the order of the pipeline, regression foundations,
least squares, maximum likelihood, $R^2$ and model adequacy. Lectures 1--4 are examined as retained
material.\\[3pt]
\textbf{Instructions.} Closed book. Write your answers on this paper, in the space provided. A
calculator is not needed --- the arithmetic is deliberately short. Answer every question --- there
is no negative marking. In Part C a reason is \textbf{optional}, but a correct reason can rescue a
wrong answer. In Parts D and E the marks are for the reasoning, not the vocabulary; an answer that
states what it cannot determine earns credit, one that asserts without a reason does not.\\[3pt]
\textbf{About the set letter.} Sets are handed out \textbf{at random}. Every set covers the same
topics at the same level, for the same marks, against the same marking scheme.}}
\vspace{3pt}


\parthead{A --- Multiple choice}{5 questions $\times$ 1 = 5 marks. Circle one letter.}

\begin{qA}
  \item \emph{A scaler fitted inside each fold gives an honest score.} An imputer's median must likewise be computed from\\
    (a) the training rows only \quad (b) all rows, before the split\\
    (c) the test rows \quad (d) the complete rows only
  \item \emph{Nearest neighbours measures distance, so it needs scaling.} A decision tree, which splits one column at a time, is therefore\\
    (a) more disturbed by scaling than neighbours are \quad (b) equally disturbed\\
    (c) unaffected by scaling \quad (d) impossible to fit on unscaled data
  \item \emph{The first principal component is the direction of largest variance.} The second is the direction of largest variance\\
    (a) among the original features only \quad (b) among the rows, not the columns\\
    (c) with the largest mean \quad (d) among those at right angles to the first
  \item \emph{A loss that falls and flattens is healthy.} A loss that rises and becomes \texttt{nan} means\\
    (a) the batch size is too small \quad (b) $\eta$ is past $2/L$\\
    (c) the model is too small \quad (d) the classes are imbalanced
  \item \emph{Predicting a labelled continuous target is supervised regression.} Sorting unlabelled customers into groups is\\
    (a) supervised regression \quad (b) supervised classification\\
    (c) unsupervised learning \quad (d) reinforcement learning
\end{qA}

\vspace{4pt}

\parthead{B --- Fill in the blanks}{4 $\times$ 1 = 4 marks. Write the missing word on the rule.}

\begin{qB}

  \item \emph{$b_0 = \bar y - b_1\bar x$ once the slope is known.} The slope itself is $b_1 = S_{xy}/\BL$.

  \item \emph{$\mathrm{SSR} = b_1^2 S_{xx}$.} And the three sums are tied together by $\mathrm{SST} = \mathrm{SSR} + \BL$.

  \item \emph{With $50$ ill patients in $1000$, always saying ``not ill'' gets $950$ right.} That is an accuracy of \BL.

  \item \emph{Supervised learning has labels; unsupervised has none.} The third family is \BL.

\end{qB}

\vfill
\begin{center}\small
\begin{tabular}{|>{\centering\arraybackslash}m{1.9cm}|*{5}{>{\centering\arraybackslash}m{1.6cm}|}}
\hline
\rule{0pt}{12pt}\textbf{Question} & A1 & A2 & A3 & A4 & A5\\[2pt]\hline
\rule{0pt}{20pt}\textbf{Answer} & & & & & \\[7pt]\hline
\end{tabular}\\[2pt]
{\footnotesize Copy your circled Part A letters into this grid. The grid is what is marked.}
\end{center}

\newpage

\parthead{C --- True or false}{3 $\times$ 1 = 3 marks. Circle one. A reason is optional, but write one --- it can rescue a wrong answer.}

\begin{qC}

  \item \emph{A \texttt{Pipeline} stops a scaler seeing the test rows.} It therefore also stops a feature computed from the target being used. \hfill \textbf{TRUE} \; / \; \textbf{FALSE}\\[1pt]
    {\small Reason:}\,\blk{140mm}\rules{1}{6.4mm}

  \item \emph{Training $R^2$ lies between $0$ and $1$.} Held-out $R^2$ is therefore also bounded below by $0$. \hfill \textbf{TRUE} \; / \; \textbf{FALSE}\\[1pt]
    {\small Reason:}\,\blk{140mm}\rules{1}{6.4mm}

  \item \emph{Cross-entropy is smooth, so gradient descent can minimise it.} Accuracy is a step function, so it cannot. \hfill \textbf{TRUE} \; / \; \textbf{FALSE}\\[1pt]
    {\small Reason:}\,\blk{140mm}\rules{1}{6.4mm}

\end{qC}

\vspace{4pt}

\parthead{D --- Short answer}{2 questions $\times$ 2 = 4 marks. Show how you decide, not just what you decide.}

\begin{qD}

  \item \emph{For $(0,1)$ and $(2,5)$ the slope is $S_{xy}/S_{xx} = 4/2 = 2$ and the intercept is $3 - 2(1) = 1$.} Now fit $\hat y = b_0 + b_1x$ to $(1,18)$, $(3,10)$, $(5,4)$, $(7,0)$ by the same route.
    \begin{enumerate}[label=(\alph*),leftmargin=7mm,itemsep=1pt,topsep=2pt]
      \item Compute $b_1$ and $b_0$, showing $S_{xx}$ and $S_{xy}$.
      \item Without working out every residual, state what $\sum_i e_i$ must come to, and say why.
    \end{enumerate}\par\vspace{62mm}

\end{qD}

\newpage

\parthead{D --- Short answer, continued}{}

\begin{qD}[start=2]

  \item \emph{A loss that falls and then flattens onto a level has converged.} A loss that falls and then bounces up and down about a level, never settling, has not.
    \begin{enumerate}[label=(\alph*),leftmargin=7mm,itemsep=1pt,topsep=2pt]
      \item Give \emph{two} distinct causes that would both produce this.
      \item State one thing you could change that would tell the two apart, and say what each outcome would mean.
    \end{enumerate}\par\vspace{70mm}

\end{qD}

\vspace{4pt}

\parthead{E --- Reasoning}{1 question $\times$ 4 = 4 marks. You are marked on the reasoning you make visible.}

\begin{qE}

  \item \emph{Worked analogue. A scaler fitted before the split leaks, because the test rows helped compute the
mean. Moving it after the split fixes it. But a column that is secretly the answer leaks no matter where it
sits --- moving it changes nothing.}

\vspace{2mm}
Now this case. The six steps below are in no particular order, and one feature, \texttt{price\_per\_sqft},
was computed as $\text{Price}/\text{Area}$, where $\text{Price}$ is the target.

\vspace{2mm}
\hspace*{4mm}\textbf{(1)} Standardise the numeric columns. \quad \textbf{(2)} Split into training and test sets.\\
\hspace*{4mm}\textbf{(3)} Median-impute the missing areas. \quad \textbf{(4)} Drop the \texttt{ID} column.\\
\hspace*{4mm}\textbf{(5)} Oversample the minority class. \quad \textbf{(6)} Fit the model.


    \vspace{2mm}
    Answer all four parts overleaf.
    \begin{enumerate}[label=(\roman*),leftmargin=7mm,itemsep=2pt,topsep=2pt]
      \item Put the six steps into a safe order. Numbers alone are enough. \hfill \textbf{[1]}
      \item One step, placed wrongly, contaminates every fitted step after it. Name it and say why.
            \hfill \textbf{[1]}
      \item One of the problems here would survive even a perfectly built \texttt{Pipeline} with the
            steps in your order. Which, and why? \hfill \textbf{[1]}
      \item One step on the list is safe to carry out \emph{before} the split. Which, and what makes
            it different from the others? \hfill \textbf{[1]}
    \end{enumerate}

\end{qE}

\newpage

\parthead{E --- Reasoning, answer space}{Label each part (i), (ii), (iii). A sketch, a table or a list is as acceptable as prose.}

\vfill

{\footnotesize\itshape Continue on a continuation sheet if you need more room; label anything you write there with its question number.}

\newpage

\end{document}